% ── methods/gen_ar.tex ────────────────────────────────────────────────────────
\subsection{Conditional Generation: Autoregressive Dual-Pathway Decoder}
\label{sec:methods_gen}

\paragraph{Architecture.}
The autoregressive (AR) generator (\textbf{v4}) attaches a frozen JEPA
encoder $f_\theta$ to a lightweight bottleneck adapter $A_\psi$ and an
autoregressive decoder with \emph{dual-pathway} conditioning
(Figure~\ref{fig:arch}C):
\begin{enumerate}[noitemsep]
  \item \textbf{Cross-attention pathway}: the condition vector $\mathbf{c}$
        is projected to a memory token prepended to the encoder output,
        so every decoder layer can attend to the condition.
  \item \textbf{AdaLN pathway}: per-layer adaptive LayerNorm scale and shift
        $(\boldsymbol{\gamma}_l, \boldsymbol{\beta}_l) =
        \text{Linear}(\mathbf{c}_\text{emb})$
        modulate each of the three sub-blocks (self-attention, cross-attention,
        FFN) independently.
\end{enumerate}
AdaLN weights are zero-initialised so training starts from an unconditioned
model.
The decoder has 4 layers ($d=384$, 8 heads, $d_\text{ff}=1536$),
yielding 10.7M trainable parameters alongside 14.2M frozen encoder parameters.

\paragraph{Condition encoding.}
Sequences are conditioned on normalised length, net charge, and
Kyte--Doolittle hydrophobicity (GRAVY):
\begin{equation}
\mathbf{c} = \bigl[\ell/50,\;\tanh(q/5),\;\tanh(h)\bigr] \in \mathbb{R}^3,
\end{equation}
projected to $\mathbf{c}_\text{emb} \in \mathbb{R}^{128}$ by a two-layer MLP.

\paragraph{Training regularisation.}
Condition dropout ($p=0.30$) and context dropout ($p=0.15$) discourage the
decoder from ignoring the condition or relying entirely on the context prefix.
The training loss combines cross-entropy token prediction with a differentiable
auxiliary loss on expected charge and GRAVY:
\begin{equation}
  \mathcal{L} = \mathcal{L}_\text{CE}
  + 0.3\,\mathcal{L}_\text{phys},
  \label{eq:gen_loss}
\end{equation}
where $\mathcal{L}_\text{phys}$ is the sum of MSE between expected (soft)
charge/GRAVY and the target values.

\paragraph{Decoder component ablations.}
To attribute charge-control performance to specific design choices, we train
two ablation variants on the same corpus and evaluate with the same
formal generation-control protocol:
\begin{itemize}[noitemsep]
  \item \textbf{v4-no-aux}: $\alpha=0$ in Eq.~\ref{eq:gen_loss}
        (no auxiliary physicochemical loss).
  \item \textbf{v4-no-dropout}: $p_\text{cond}=p_\text{ctx}=0$
        (no CFG-style dropout).
\end{itemize}
Results for these ablations are reported in Table~\ref{tab:gen_ablation}.
